\documentclass{article}
\usepackage{geometry}
\geometry{a4paper, margin=1in}
\usepackage{xcolor}
\usepackage{hyperref}
\hypersetup{colorlinks=true, linkcolor=blue, urlcolor=blue, citecolor=blue}

\title{Critical Analysis of Topic~5:\\
\textit{Self-Supervised Representation Learning for Morphogenetic Trajectory Discovery from Light-Sheet Microscopy}}
\author{Devil's Advocate Review}
\date{\today}

\begin{document}
\maketitle

\noindent This document is a self-critique of the proposal in \texttt{topic5.tex}. It
identifies three flaws of decreasing severity---one fundamental, one significant,
one moderate---and proposes concrete fixes for each. The goal is to stress-test
the idea \emph{before} presenting it to potential supervisors.

%% =========================================================================
\section{Flaw 1 --- The $n{=}1$ Trajectory Problem (FUNDAMENTAL)}
%% =========================================================================

\subsection*{The claim in the proposal}
``We apply trajectory inference methods adapted from the single-cell genomics
literature---specifically optimal transport (Waddington-OT) or diffusion
maps---to reconstruct continuous morphogenetic trajectories through the learned
embedding space.''

\subsection*{Why it cannot work as stated}
Waddington-OT~\cite{wot} was designed for a very specific data regime:

\begin{enumerate}
	\item At each timepoint, \textbf{thousands of cells} are measured (a whole
	      population snapshot).
	\item Measurement is \textbf{destructive}---cells are lysed for RNA
	      extraction, so no cell can appear at two timepoints.
	\item OT infers the \textbf{missing temporal couplings} between these
	      population snapshots, recovering probable ancestor--descendant
	      relationships across timepoints that were never directly observed.
\end{enumerate}

\noindent In a time-lapse video of a \textbf{single embryo}, none of these
conditions hold:

\begin{itemize}
	\item There is exactly \textbf{one specimen} per timepoint, not a population.
	\item The ``trajectory'' is trivially the \textbf{temporal sequence of
		      frames}---there is nothing to infer. The coupling between frame~$t$ and
	      frame~$t{+}1$ is the identity: it is the same embryo one timestep later.
	\item Optimal transport between two point masses (single-point distributions)
	      degenerates to a trivial map. The Sinkhorn solver has nothing
	      meaningful to compute.
\end{itemize}

\noindent This is not a minor technical issue but a \textbf{category error}: the
proposal applies a method designed for population-snapshot data to
continuous single-specimen time-lapse data, where the problem the method solves
does not exist.

\subsection*{Supporting literature}

\begin{itemize}
	\item Wang et al.\ (2020,
	      \href{https://doi.org/10.1016/j.cels.2020.11.006}{\textit{Cell
			      Systems}}) explicitly argue that time-lapse imaging is \emph{more
		      informative} than what OT reconstructs from snapshots, precisely because
	      the temporal identity of each object is preserved across frames.
	\item Copperman et al.\ (2023,
	      \href{https://doi.org/10.1073/pnas.2216437120}{\textit{PNAS}}) analyse
	      live-cell imaging trajectories using trajectory embedding and diffusion
	      analysis \emph{instead of} OT, because the continuous time-lapse
	      context makes OT redundant.
	\item No published work applies Waddington-OT to single-specimen time-lapse
	      imaging data---this absence is not accidental.
\end{itemize}

\subsection*{Proposed fix}

Two options, either of which rescues the proposal:

\begin{description}
	\item[Option A --- Drop OT, reframe as time-series analysis.]
	      Replace Waddington-OT with methods suited to single-trajectory temporal
	      data: HMMs, change-point detection, time-series segmentation, or the
	      diffusion-map approach of Copperman et al.\ (2023). The core
	      contribution becomes: \emph{learn a latent space via TAP + DINOv2, then
		      segment the single developmental trajectory into unsupervised stages.}

	\item[Option B --- Pivot to multi-embryo population analysis.]
	      The \textbf{SLICE-1} dataset (Strobl et al., \textit{Scientific Data},
	      Dec 2025~\cite{slice1}) provides \textbf{50 Tribolium embryos} (5.7\,M
	      images, 2.7\,TB, five transgenic lines with triplicates). Kraemer et
	      al.\ (bioRxiv, Jan 2026~\cite{kraemer2026}) contribute 16 additional
	      isotropic 3D Tribolium gastrulation datasets. With 50+ embryos, one
	      \emph{can} treat each timepoint as a population snapshot across specimens
	      and apply OT to infer how the \emph{distribution} of developmental
	      states evolves. This makes OT valid and scientifically interesting
	      (e.g., comparing wild-type vs.\ mutant trajectories).
\end{description}

%% =========================================================================
\section{Flaw 2 --- DINOv2 Domain Gap (SIGNIFICANT)}
%% =========================================================================

\subsection*{The claim in the proposal}
``Spatial representations extracted from a frozen DINOv2 vision foundation
model [\ldots] providing rich per-frame feature embeddings without any labels.''

\subsection*{Why it is problematic}
DINOv2~\cite{dinov2} was trained on 142\,M curated natural RGB photographs
(LVD-142M). Tribolium light-sheet cartographic projections are:

\begin{itemize}
	\item \textbf{Grayscale}, typically 16-bit.
	\item \textbf{Sparse fluorescence signal} on a dark background---visually
	      nothing like natural photographs.
	\item \textbf{Cartographic projections} of 3D volumes, introducing geometric
	      distortions absent from natural images.
\end{itemize}

\noindent Frozen DINOv2 features will not be \emph{garbage}---the model still
captures generic texture and spatial statistics---but they will be
\textbf{substantially degraded} compared to what a domain-aware model would
produce.

\subsection*{Supporting literature}

\begin{itemize}
	\item \textbf{DINOv3 medical benchmark} (Liu et al., 2026): explicitly
	      identifies performance degradation on electron microscopy and
	      specialized non-photographic modalities.
	\item \textbf{Cell-DINO} (PLOS Computational Biology, 2025): domain-specific
	      retraining of DINO on fluorescence microscopy improves downstream
	      performance by $>20\%$ over frozen DINOv2 features.
	\item \textbf{ViTally Consistent} (Recursion Pharmaceuticals, ICML 2025):
	      domain-specific MAE pretraining on high-content screening images yields
	      $\sim$60\% improvement over natural-image baselines on biological
	      retrieval tasks.
\end{itemize}

\subsection*{Proposed fix}

\begin{itemize}
	\item Do \textbf{not} assume frozen DINOv2 features work out of the box.
	      Plan for at least \textbf{LoRA fine-tuning} (low-rank adaptation) of
	      DINOv2 on the target microscopy data, which adds minimal parameters
	      and is feasible within a one-semester scope.
	\item Alternatively, extract features from \textbf{intermediate layers}
	      rather than only the final CLS token; intermediate layers often retain
	      more transferable spatial information for out-of-distribution inputs.
	\item Treat ``How well do frozen foundation-model features transfer to
	      light-sheet data?'' as an \textbf{explicit experimental question}
	      with a planned comparison (frozen vs.\ LoRA-adapted vs.\ MAE
	      pretrained on microscopy), rather than an untested assumption.
\end{itemize}

%% =========================================================================
\section{Flaw 3 --- TAP Untested at Tissue / Embryo Scale (MODERATE)}
%% =========================================================================

\subsection*{The claim in the proposal}
TAP ``captures biologically meaningful events (e.g., cell division) without any
labels'' and is applied directly to whole-embryo light-sheet data.

\subsection*{Why it is uncertain}
TAP~\cite{tap} has been validated exclusively on \textbf{2D cell culture
	videos}, where the time-arrow signal comes from sharp, spatially localized
events: individual cell divisions, migrations, and apoptosis. These events
produce clear frame-to-frame asymmetry that a convolutional network can detect.

Whole-embryo morphogenesis operates at a different spatial scale:

\begin{itemize}
	\item Key events (gastrulation, germband elongation, dorsal closure) are
	      \textbf{global tissue deformations}, not localized pixel-level
	      discontinuities.
	\item Temporal changes may appear ``smooth'' at the frame-to-frame level,
	      reducing the arrow-of-time signal that TAP relies on.
	\item The original TAP paper~\cite{tap} does not test on any tissue-scale or
	      embryo-scale data.
\end{itemize}

\noindent However, embryonic development \emph{is} thermodynamically
irreversible (entropy production, symmetry breaking, irreversible cell fate
commitment), so a time-arrow signal almost certainly exists---the question is
whether a standard convolutional TAP architecture can detect it at this scale.

\subsection*{Proposed fix}

\begin{itemize}
	\item Frame TAP-on-embryo-data as an \textbf{open research question and
		      genuine contribution}, not as an assumption. If it works, that is a
	      novel finding. If it fails, diagnosing \emph{why} (resolution,
	      temporal sampling rate, architecture) is itself valuable.
	\item Plan a \textbf{multi-scale TAP architecture}: apply TAP at both
	      local patches (where cell-level events are visible) and global scale
	      (whole-embryo), then compare which scale carries more temporal signal.
	\item Include a \textbf{fallback self-supervised signal}: if TAP proves
	      weak at embryo scale, use contrastive temporal objectives (e.g.,
	      temporal contrastive learning with positive pairs from nearby frames)
	      that do not require a crisp arrow-of-time signal.
\end{itemize}

%% =========================================================================
\section{Summary}
%% =========================================================================

\begin{center}
	\begin{tabular}{|l|l|p{7.5cm}|}
		\hline
		\textbf{Flaw}                          & \textbf{Severity} & \textbf{Verdict} \\
		\hline
		$n{=}1$ trajectory / OT misapplication & Fundamental       &
		\textbf{Must be fixed} before presenting to supervisors. Either drop OT or
		pivot to multi-embryo (SLICE-1).                                              \\
		\hline
		DINOv2 domain gap                      & Significant       &
		Plan for LoRA fine-tuning; make transfer quality an experimental variable,
		not an assumption.                                                            \\
		\hline
		TAP at tissue scale                    & Moderate          &
		Frame as a research question, not a given. Plan fallback self-supervised
		objectives.                                                                   \\
		\hline
	\end{tabular}
\end{center}

\noindent\textbf{Bottom line:} The core idea---self-supervised representations
for unsupervised developmental trajectory analysis---is sound and timely. But
the specific pipeline described in \texttt{topic5.tex} will not work as written
due to the OT category error. With the fixes above (especially Option~B using
SLICE-1), the proposal becomes both technically correct and more ambitious.

%% =========================================================================
\begin{thebibliography}{9}

	\bibitem{wot} G.~Schiebinger et al.
	``Optimal-transport analysis of single-cell gene expression identifies
	developmental trajectories in reprogramming.''
	\textit{Cell}, 176(4):928--943, 2019.

	\bibitem{tap} B.~Gallusser, M.~Stieber, and M.~Weigert.
	``Self-supervised dense representation learning for live-cell microscopy
	with time arrow prediction.''
	\textit{MICCAI}, 2023.

	\bibitem{dinov2} M.~Oquab et al.
	``DINOv2: Learning robust visual features without supervision.''
	\textit{TMLR}, 2024.

	\bibitem{slice1} F.~Strobl et al.
	``SLICE-1: A comprehensive light-sheet imaging dataset of \textit{Tribolium
		castaneum} embryogenesis.''
	\textit{Scientific Data}, Dec 2025.

	\bibitem{kraemer2026} E.~T.~Kraemer et al.
	``Isotropic reconstruction of \textit{Tribolium castaneum} gastrulation
	from multiview light-sheet microscopy.''
	\textit{bioRxiv}, Jan 2026.

\end{thebibliography}

\end{document}
